\documentclass[paper=a4, fontsize=11pt]{jlreq} % 日本語用の標準的なクラス

% --- パッケージ群 ---
\usepackage[utf8]{inputenc}
\usepackage[T1]{fontenc}
\usepackage{amsmath, amssymb, amsthm} % 数学の基本
\usepackage{physics} % \ket{0}, \pdv{x} などが簡単に打てる
\usepackage{bm} % 太字ベクトル \bm{v}
\usepackage{graphicx}
\usepackage{hyperref} % リンク・目次用
\usepackage{cite} % 引用用

\newtheoremstyle{break}% スタイル名
  {}% 上部のスペース（デフォルト）
  {}% 下部のスペース（デフォルト）
  {\normalfont}% 本文のフォント（日本語なので斜体にしない）
  {}% インデント量
  {\bfseries}% 見出しのフォント（太字）
  {}% 見出し後の句読点（不要なら空に）
  {\newline}% ★ここで改行を指定！★
  {}% 見出しのカスタム指定（デフォルト）

% --- 定理環境の設定 ---
\theoremstyle{break}
\newtheorem{theorem}{定理}[section]
\newtheorem{definition}[theorem]{定義}
\newtheorem{proposition}[theorem]{命題}
\newtheorem{example}[theorem]{例}
\newtheorem{problem}[theorem]{問題} % 自作問題用
\newtheorem*{fact}{事実}
\newtheorem*{setup}{設定}
\newtheorem*{connect}{$\mho\mho$}

% --- 番号なしの環境（地の文の代わり） ---
\theoremstyle{remark} % 見た目を少し控えめにするスタイル
\newtheorem*{remark}{※}

% --- 独自のショートカット (CFT/パンルヴェ用) ---
\newcommand{\Virexp}[2]{[L_{#1}, L_{#2}]} % ヴィラソロの交換関係
\newcommand{\taufunc}{\tau(z)} % タウ関数

\title{パンルヴェ方程式に関する勉強ノート}
\author{松浦光佑}
\date{\today}

\begin{document}

\maketitle
\tableofcontents % 自動で目次を作成

\section{はじめに}
本ノートは、岡本『パンルヴェ方程式』に関する学習記録である。
日々の学習を「AC（Accepted）」として記録し、GitHubでの進捗可視化を目的とする。

\section{フックス型微分方程式}

\subsection{2階線型常微分方程式}

\begin{setup}
    \begin{equation}\label{eq:2-ode}
    \frac{d^{2}y}{dx^{2}} + p_{1}(x)\frac{dy}{dx} + p_{2}(x)y = 0   
    \end{equation}
    を考える。$x\in\mathbb{C}$とする。よって解$y$は解析的である。
    また、$p_{1}(x),p_{2}(x)$は$x$の有理関数とする。よって\eqref{eq:2-ode}は広くても$\mathbb{P}^{1}(\mathbb{C})$上の微分方程式である。
\end{setup}

\begin{definition}[$x=\infty$における微分方程式]
    \eqref{eq:2-ode}で$x=\frac{1}{u}$とする。u=0の近傍で考えたものを\textbf{$x=\infty$における微分方程式}という。ある有理関数$q_{1},q_{2}$を用いて
    \[
    \frac{d^{2}y}{du^{2}} + q_{1}(u)\frac{dy}{du} + q_{2}(u)y = 0
    \]
    のように書いたとき、\eqref{eq:2-ode}に連鎖律 $\frac{d}{dx} = -u^2 \frac{d}{du}$ を用いたものとの係数比較により、
    \[
    q_{1}(u) = \frac{2}{u}-\frac{1}{u^{2}}p_{1}\left(\frac{1}{u}\right),\, q_2(u)=\frac{1}{u^{4}}p_{2}\left(\frac{1}{u}\right)
    \]
    という関係がわかる。
    \begin{itemize}
        \item[\textbf{Insight:}] 複素幾何学（リーマン球面 $\mathbb{P}^1(\mathbb{C})$）の観点では、無限遠点 $\infty$ は有限の点と全く対等な「ただの1点」である。$u=1/x$ はその点を見るための「レンズ（局所座標）」に過ぎない。$q_1(u)$ に現れる $2/u$ という余分な項は、座標変換の非線形性から生じる見かけ上のズレであり、$x=\infty$ での特異性の度合い（確定特異点かどうか）を判定する上で極めて重要な役割を果たす。
        \item[\textbf{Link:}] 岡本2.1節のp.49を参照。
    \end{itemize}
\end{definition}

\begin{connect}
    $p_{1},p_{2}$を$x=x_{0}$で正則とする。このとき、コーシーの存在定理から、任意の$(y_{0},y'_{0})\in\mathbb{C}^{2}$に対して、
    \[
    y(x_{0}) = y_{0},\,\frac{dy}{dx}(x_{0})=y_{0}'
    \]
    を満たす解$y=\varphi(x;x_{0},y_{0},y_{0}')$がただ一つ存在する。ここで、$u=x-x_{0}$とする。テイラー展開は
    \[
    \varphi(x;x_{0},y_{0},y_{0}') = \sum_{j=0}^{\infty}c_{j}u^{j},\,c_{0}=y_{0},\,c_{1}=y_{0}'
    \]
    と表せる。また、
    \[
    p_{1}(x) = \sum_{j=0}^{\infty}a_{j}u^{j},\,p_{2}(x) = \sum_{j=0}^{\infty}b_{j}u^{j}
    \]
    と表すことにする。
\end{connect}

\begin{proposition}
    $x=\xi$を$y=\varphi(x;x_{0};y_{0},y_{0}')$の特異点とすれば、$\xi$は$p_{1}(x)$または、$p_{2}(x)$の特異点である。\\
    \textbf{証明のスケッチ}:
    証明は一旦略
    \begin{itemize}
        \item[\textbf{Insight:}] 線形微分方程式の最大の特長は、「解の特異点が初期条件 $(x_{0}, y_{0}, y_{0}')$ に一切依存しない」ことである。特異点はあらかじめ係数の特異点として空間上に「固定」されており、どこから出発しようと障害物の位置は変わらない。これは、非線形方程式における「動く特異点」の異常さを際立たせるための、最も重要な比較基準（ベースライン）となる。
        \item[\textbf{Link:}] 岡本2.1節の命題2.1/p.50を参照。
    \end{itemize}
\end{proposition}


\subsection{確定特異点}

    $x=\xi$をフックス型微分方程式
    \begin{equation}\label{eq:Fuchs}
    \frac{d^{2}y}{dx^{2}} + p_{1}(x)\frac{dy}{dx} + p_{2}(x)y = 0   
    \end{equation}
    の特異点とする。

\begin{proposition}[]
    % --- ここに数式 ---
    適当な2次の行列関数$\Phi_{\xi}(x)$を選んで
    \[
    \phi(x) = \vec{z}(x)\Phi_{\xi}(x)
    \]
    で定まる$\vec{z}(x)$が$x=\xi$のまわりで1価になるようにできる。もし、$\Gamma$が対角化可能ならば、$\Phi_{\xi}(x)$は対角行列になるように選べる。\\
    また、1価性により$x=\xi$のまわりで
    \begin{equation}\label{eq:power_series_z}
    \vec{z}(x) = \sum_{j=-\infty}^{\infty}\vec{a_{j}}u^{j}
    \end{equation}
    と表せる。
    \textbf{証明のスケッチ:}
    
    \begin{itemize}
        \item[\textbf{Insight:}] 通常、特性指数の差が整数の場合、解の片方に $\log$（対数項）が現れる。しかし、係数が特定の条件を満たして「$\log$ の項が完全に消滅する」特殊なケースがあり、それが非対数的特異点である。「みかけの特異点」の周りで解が分岐しないためにはこの性質が必須であり、この「対数項が消える条件」を計算することが、後のハミルトニアン導出の鍵を握る。
        \item[\textbf{Link:}] 岡本2.2節の命題2.5/p.59を参照。
    \end{itemize}
\end{proposition}

\begin{definition}[微分方程式の確定特異点]
    % --- ここに数式 ---
    \eqref{eq:power_series_z}において、ある整数$j_{0}$があって、$j\leq j_{0}$となるすべてのjに対し、
    \[
    \vec{a_{j}} = \vec{0}
    \]
    となる時$\xi$は\textbf{確定特異点}という。確定特異点でない時\textbf{不確定特異点}という。
    \begin{itemize}
        \item[\textbf{Insight:}] 特異点周りの級数展開において、「負のべき乗の項が無限に続かず、どこかで打ち止めになる（有限個しかない）」状態を表している。複素関数論における「極（pole）」に相当する、解析的に扱いやすい「おとなしい特異点」の定義である。パンルヴェ方程式の根幹である「動く特異点が極のみである（パンルヴェ性）」という性質を語る上で、この「確定（極）か、不確定（真性特異点など）か」の厳密な線引きがすべての土台となる。
        \item[\textbf{Link:}] 岡本2.2節の定義2.4/p.59を参照。
    \end{itemize}
\end{definition}

    $\theta_{1},\theta_{2}\in\mathbb{R}$と$r>0$に対して、
    \[
    S(\theta_{1},\theta_{2}) = \{u\,|\, 0<|x|<r, \theta_{1} < \arg u < \theta_{2} \}
    \]
    とおく。ここで$u=x-\xi$である。

\begin{definition}[解析的な解の確定特異点]
    % --- ここに数式 ---
    $x=\xi$で特異点をもつ解析関数の分枝$\varphi(x)$に対して、$S(\theta_{1},\theta_{2})$で$\varphi(x)$が定義されるように$r$を十分小さくとる。
    ある$A>0$がとれて$\theta_{1}<\theta_{2}$に対し、
    \[
    \lim_{u\to 0, u\in S(\theta_{1},\theta_{2})}|u|^{A}|\varphi(x)| = 0
    \]
    となるとき、$\varphi(x)$は$x=\xi$に\textbf{確定特異点}を持つという。
    \begin{itemize}
        \item[\textbf{Insight:}] 方程式の解が対数関数などの多価性を持つため、特異点への近づき方を「特定の扇形領域（セクター）内」に制限して評価しているのがポイント。条件式は「特異点に近づく際の発散スピードが、せいぜい多項式オーダー（$|u|^{-A}$ 以下）に収まり、指数関数的（$e^{1/|u|}$ など）には爆発しない」ことを意味する。代数的な「おとなしさ（高々極）」を、解析的な「増大度のバウンド」として厳密に定式化したものである。
        \item[\textbf{Link:}] 岡本2.2節の定義2.5/p.60を参照。
    \end{itemize}
\end{definition}

\begin{remark}[特異点 $\omega$ とベース空間の座標 $x=\xi$ の区別について]
    原著において、解析関数 $\varphi(x)$ の特異点を表す際、ベースとなる複素平面上の座標 $x=\xi$ とは別に $\omega$ という記号が用いられることがある。
    これは、$\varphi(x)$ が対数関数や累乗根のような多価関数（特異点の周りを一周すると値が変わる関数）の分枝（branch）であることを考慮した厳密な表記である。
    
    複素幾何学（リーマン面の理論）の立場では、ベース空間の $x=\xi$ という1点の上にも、分枝の選び方によってリーマン面上の異なる複数の点 $\omega_1, \omega_2, \dots$ が重なって存在し得るため、これらを区別している。
    \begin{itemize}
        \item[\textbf{Insight:}] 確定特異点の定義における極限評価 $\lim |u|^A|\varphi(x)|=0$ では、ベース空間での距離 $|u| = |x-\xi|$ のみが計算に関与する。したがって、「特定の分枝を選び、扇形領域内で近づく」という前提さえあれば、リーマン面上の点 $\omega$ を明示しなくても解析的な論理は完結する。無駄を省き本質にフォーカスするため、本ノートの数式展開においては原則として $x=\xi$ のみに表記を統一する。
    \end{itemize}
\end{remark}

\begin{definition}[非対数的特異点]
    % --- ここに数式 ---
    $x=\xi$の特性指数$s_{1},s_{2}$が整数差のとき$x=\xi$は\textbf{非対数的特異点}と呼ばれる。
    \begin{itemize}
        \item[\textbf{Insight:}] 通常、特性指数の差が整数の場合、解の片方に $\log$（対数項）が現れる。しかし、係数が特定の条件を満たして「$\log$ の項が完全に消滅する」特殊なケースがあり、それが非対数的特異点である。「みかけの特異点」の周りで解が分岐しないためにはこの性質が必須であり、この「対数項が消える条件」を計算することが、後のハミルトニアン導出の鍵を握る。
        \item[\textbf{Link:}] 岡本2.2節の(2.20)/p.64を参照。
    \end{itemize}
\end{definition}

\subsection{不確定特異点}

\subsection{高階フックス型線型常微分方程式}

\begin{definition}[リーマン図式]
    % --- ここに数式 ---
    \[
    \left\{
    \begin{matrix}
        x = \xi_1 & x = \xi_2 & \cdots & x = \xi_m & x = \infty \\
        s_1^{(1)} & s_1^{(2)} & \cdots & s_1^{(m)} & s_1^{(\infty)} \\
        s_2^{(1)} & s_2^{(2)} & \cdots & s_2^{(m)} & s_2^{(\infty)} \\
        \vdots    & \vdots    &        & \vdots    & \vdots \\
        s_n^{(1)} & s_n^{(2)} & \cdots & s_n^{(m)} & s_n^{(\infty)}
    \end{matrix}
    \right\}
    \]
    \begin{itemize}
        \item[\textbf{Type:}] 確定特異点$\xi_{i}$と各確定特異点の特性指数$s_j^{(i)}$を書き並べたもの。
        \item[\textbf{Insight:}] この表の数値（局所的な性質）を\textbf{一切変えずに}、特異点の位置 $\xi_i$ だけを時間発展させた時に、裏で満たされるべき非線形の方程式が「ガルニエ系」。つまり、モノドロミー保存変形において\textbf{「絶対に定数として保存したいターゲット」}がこれ。
        \item[\textbf{Link:}] 岡本2.4節の(2.48)/p.81を参照。
    \end{itemize}
\end{definition}

\subsection{みかけの特異点}

\begin{definition}[みかけの特異点]
  % --- ここに数式 ---
    $\Xi$の元を確定特異点としてもつ微分方程式を考える。
    このときモノドロミー$\hat{\rho}$が等しく、さらに$\Xi\cup\Lambda$の元を確定特異点としてもつフックス型微分方程式が存在したとする。
    このとき、以下の2条件を満たす$\Lambda$の各元を\textbf{みかけの特異点}という。
    \begin{itemize}
        \item[(1)] $x=\lambda$における特性指数は全て整数。
        \item[(2)] $x=\lambda$は非対数的特異点
    \end{itemize}
    \begin{itemize}
        \item[\textbf{Insight:}] 方程式の係数には特異性が現れるが、解自体はそこで分岐を持たず正則（局所モノドロミーが単位行列）になる「偽物の特異点」。パンルヴェ方程式やガルニエ系の理論において最大のポイントは、**「このみかけの特異点の位置 $\lambda$ が、ハミルトン系における正準変数（未知関数 $q$）の役割を果たす」**という点にある。
        \item[\textbf{Link:}] 岡本2.5節/p.86を参照。
    \end{itemize}
\end{definition}

\subsection{連立微分方程式系と単独高階微分方程式}

\subsection{リーマンの問題}

\subsection{フックスの問題}

\begin{definition}[モノドロミー保存変形]
  % --- ここに数式 ---
    \begin{equation}\label{eq:before}
    \frac{d^{n}y}{dx^{n}}+\sum_{j=1}^{n} p_{j}(x)\frac{d^{n-1}y}{dx^{n-1}} = 0   
    \end{equation}
    という以下の(P1)-(P3)を満たすフックス型微分方程式を考える。
    \begin{itemize}
        \item[(P1)] $x=\xi_{1},...,\xi_{n},\xi_{n+1}=\infty$を確定特異点、$x=\lambda_{1},...,\lambda_{g}$をみかけの特異点としてもつ。
        \item[(P2)] 各$x=\xi_{i}$における特性指数のどの2つも、差が整数ではない。
        \item[(P3)] モノドロミー$\hat{\rho}$は既約である。
    \end{itemize}
    $p_{j}(x)$が$t$に解析的に依存しているとする。
    \begin{equation}\label{eq:after}
    \frac{d^{n}y}{dx^{n}}+\sum_{j=1}^{n} p_{j}(x,t)\frac{d^{n-1}y}{dx^{n-1}} = 0   
    \end{equation}
    が以下の条件(H)を満たす時、\eqref{eq:after}を\eqref{eq:before}の\textbf{モノドロミー保存変形}または\textbf{ホロノミック変形}という。
    \begin{itemize}
        \item[\textbf{Type:}] \eqref{eq:before}は$\mathbb{P}^{1}(\mathbb{C})$上定義されている。$U\subset\mathbb{C}^{L}$は領域。$t=(t_{1},...,t_{L})\in U$
        \item[\textbf{Insight:}] 方程式のパラメータ $t$ を動かしても、基本群の表現（モノドロミー）が変化しないような変形。Schlesinger系やPainlevé方程式の導出における出発点となる。
        \item[\textbf{Link:}] 岡本2.8節の定義2.14/p.105を参照。
    \end{itemize}
\end{definition}

\appendix
\section{補遺：定理・命題の厳密な証明}

\subsection{命題2.2（線形方程式の固定特異点）の証明}
% ここに、後日時間があるときに「解析接続と収束半径の下限」を用いた
% 厳密な証明をUpsolveとして記述する。
% （※すでに完成している証明コードを貼り付けてもよい）

$x = \xi$ が解 $y = \varphi(x; x_0, y_0, y'_0)$ の特異点であるとする。
    このとき、複素解析的な解の特異点の定義より、初期値 $x_0$ から $\xi$ に向かうある経路 $\gamma$ が存在し、「$\gamma$ 上の $\xi$ を除くすべての点（$\gamma \setminus \{\xi\}$）において解は正則に解析接続可能である」ことが保証されている。
    解が正則である経路上では、係数 $p_1(x), p_2(x)$ の特異点を避けている（あるいは迂回できる）ため、$\gamma \setminus \{\xi\}$ 上で $p_1, p_2$ は正則であるとしてよい。

    ここで、終点 $\xi$ も $p_1(x)$ および $p_2(x)$ の正則点であると仮定する（背理法の仮定）。
    すると、端点を含む経路 $\gamma$ 全体（有界閉集合）において、$p_1, p_2$ はすべて正則となる。
    したがって、$\gamma$ 上の各点から $p_1, p_2$ の最も近い特異点（存在すれば）までの距離の下限は、ある正の定数 $\delta > 0$ 以上となる。

    線形常微分方程式の性質により、$\gamma$ 上の任意の点を中心とする解のテイラー展開の収束半径は、少なくとも $\delta$ 以上であることが保証される。
    ゆえに、$x_0$ を起点として半径 $\delta$ の収束円を有限回繋ぎ合わせる解析接続を行うことで、解は $\gamma$ に沿って特異性にぶつかることなく、確実に $\xi$ の近傍まで延長される。

    結果として、解 $y = \varphi(x)$ は $\xi$ の近傍において正則（テイラー展開可能）となる。これは「$\xi$ が解の特異点である」という前提に矛盾する。
    したがって、最初の仮定は誤りであり、$x = \xi$ は必ず $p_1(x)$ または $p_2(x)$ の特異点でなければならない。\hfill $\blacksquare$




% --- 今後の課題（TODOリスト） ---
\section{Next Step / To-Do}
\begin{itemize}
    \item [-] 確定特異点とは？
    \item [-] 特性指数とは？
    \item [-] モノドロミーが既約とはどういうこと？
    \item [-] モノドロミーとは?
\end{itemize}

\end{document}